\section{Orbit congruence and the two-affine-prime reduction}
\label{sec:tpc57-orbit-reduction}

We retain the literal hypotheses and notation of
\cref{sec:tpc57-factor-exchange}.  In particular \(h_0\ne0\) is fixed,
\(H=\operatorname{rad}|h_0|\), \(X\) is sufficiently large that
\(L>|h_0|\), and no solution is counted unless it belongs to the exact
retained carrier.

\subsection{Exact orbit congruence}

\begin{lemma}[The terminal gcd is fixed by the intercept]
\label{lem:tpc57-terminal-gcd}
For every terminal atom \(u=(d\ell,r,j,\gamma)\),
\begin{equation}
 (d\ell,r)=1,
 \qquad
 j\equiv-h_0(d\ell)^{-1}\pmod r,
 \qquad
 \boxed{(r,j)=(r,h_0).}
 \label{eq:tpc57-exact-orbit-congruence}
\end{equation}
Since \(r\) is squarefree, the common gcd in
\eqref{eq:tpc57-exact-orbit-congruence} divides \(H\).
\end{lemma}

\begin{proof}
The literal carrier gives \((d,r)=1\).  If the source prime
\(\ell\) divided \(r\), then both \(rp\) and \(d\ell j\) would be
divisible by \(\ell\), so the terminal equation would imply
\(\ell\mid h_0\), impossible for \(\ell\ge L>|h_0|\).  Hence
\((d\ell,r)=1\).  Reducing
\(rp-d\ell j=h_0\) modulo \(r\) gives the displayed congruence.
Multiplication by a unit modulo \(r\) preserves the gcd with \(r\), so
\((r,j)=(r,h_0)\).  The final assertion follows from squarefreeness.
\end{proof}

\begin{corollary}[Orbit-time sparsity for a fixed row and cofactor]
\label{cor:tpc57-fixed-row-orbit-sparsity}
Fix \(d,r,\ell\) satisfying the literal source and cofactor conditions.
In an orbit interval of diameter at most \(C_JJ\), the terminal equality
can hold for at most
\begin{equation}
 1+\left\lfloor\frac{C_JJ}{r}\right\rfloor
 \label{eq:tpc57-fixed-row-j-count}
\end{equation}
integer values of \(j\).  Intersecting with the declared orbit residue
class, profiles, or terminal-prime condition can only decrease this
number.
\end{corollary}

\begin{proof}
By \cref{lem:tpc57-terminal-gcd}, every possible \(j\) lies in one fixed
residue class modulo \(r\).  The elementary spacing bound in an interval
of diameter \(C_JJ\) proves the result.
\end{proof}

The congruence is stronger than the necessary statement
\((r,j)\mid h_0\): on the literal terminal support the gcd is exactly
\((r,h_0)\).  It does not assert that the unique residue class contains a
retained \(j\), or that the resulting quotient is prime.

\subsection{A simultaneous two-row orbit class}

Use the factor-exchange notation
\[
 d_1=ab,\quad d_2=ac,
 \qquad r_1=ce,\quad r_2=be,
\]
with \(a,b,c,e\) pairwise coprime.  For fixed
\((a,b,c,e,\ell_1,\ell_2)\), the two terminal equalities require
\begin{equation}
 \begin{aligned}
 j&\equiv-h_0(ab\ell_1)^{-1}\pmod {ce},\\
 j&\equiv-h_0(ac\ell_2)^{-1}\pmod {be}.
 \end{aligned}
 \label{eq:tpc57-two-orbit-congruences}
\end{equation}
All inverses exist by \cref{lem:tpc57-terminal-gcd}.

\begin{proposition}[Exact CRT compatibility and joint orbit spacing]
\label{prop:tpc57-joint-orbit-class}
The two congruences in
\eqref{eq:tpc57-two-orbit-congruences} are compatible if and only if
\begin{equation}
 \boxed{
 \frac e{(e,h_0)}\mid b\ell_1-c\ell_2.}
 \label{eq:tpc57-crt-compatibility}
\end{equation}
When compatible, they define one residue class modulo
\begin{equation}
 \operatorname{lcm}(ce,be)=bce.
 \label{eq:tpc57-joint-orbit-modulus}
\end{equation}
Consequently an orbit interval of diameter \(C_JJ\) contains at most
\begin{equation}
 1+\left\lfloor\frac{C_JJ}{bce}\right\rfloor
 \label{eq:tpc57-joint-orbit-count}
\end{equation}
simultaneously admissible integer times.  A fixed inherited orbit class
can only thin this set further.
\end{proposition}

\begin{proof}
The two residue classes must agree modulo
\((ce,be)=e\).  Multiplying their equality modulo \(e\) by the available
units shows that compatibility is equivalent to
\[
 e\mid h_0(b\ell_1-c\ell_2),
\]
which is exactly \eqref{eq:tpc57-crt-compatibility}.  The Chinese
remainder theorem then gives one class modulo the least common multiple
\(bce\), and the spacing estimate follows.
\end{proof}

This is an orbit-counting statement, not a prime-counting theorem.  Once
\(j\) is chosen, both terminal candidates are determined, but the CRT
does not prove either candidate prime.

\subsection{The determinant ladder}

Fix a coarse output \(y=(\gamma,s,j)\) and an admissible divisor
\(d\mid s\), and put \(r=s/d\).  The source and terminal primes must solve
\begin{equation}
 rp-dj\ell=h_0.
 \label{eq:tpc57-ladder-equation}
\end{equation}
Let
\begin{equation}
 g:=(r,dj)=(r,j),
 \label{eq:tpc57-ladder-gcd}
\end{equation}
where the second equality uses \((d,r)=1\).

\begin{proposition}[Complete nonproportional affine reduction]
\label{prop:tpc57-two-affine-reduction}
Equation \eqref{eq:tpc57-ladder-equation} has an integral solution if and
only if \(g\mid h_0\).  If \((\ell_0,p_0)\) is one solution, every
integral solution is uniquely
\begin{equation}
 \boxed{
 \ell(t)=\ell_0+\frac r g t,
 \qquad
 p(t)=p_0+\frac{dj}{g}t,
 \qquad t\in\mathbb Z.}
 \label{eq:tpc57-two-affine-forms}
\end{equation}
The two affine forms have nonzero cross determinant
\begin{equation}
 \frac{dj}{g}\ell_0-\frac r g p_0
 =-\frac{h_0}{g}\ne0.
 \label{eq:tpc57-affine-determinant}
\end{equation}
On the literal terminal support,
\begin{equation}
 g=(r,h_0)\mid H.
 \label{eq:tpc57-literal-ladder-gcd}
\end{equation}
\end{proposition}

\begin{proof}
The gcd criterion for the linear Diophantine equation is
\(g\mid h_0\).  Differences of two solutions satisfy
\(r\Delta p=dj\Delta\ell\).  After division by \(g\), the two remaining
coefficients are coprime, so
\(\Delta\ell=(r/g)t\) and
\(\Delta p=(dj/g)t\), proving
\eqref{eq:tpc57-two-affine-forms}.  Substitution of the base solution
gives \eqref{eq:tpc57-affine-determinant}.  Finally,
\cref{lem:tpc57-terminal-gcd} gives \(g=(r,j)=(r,h_0)\), and
squarefreeness of \(r\) gives the divisibility by \(H\).
\end{proof}

Let \(\mathfrak T_{\gamma,s,j,d}\subset\mathbb Z\) be the exact set of
parameters \(t\) in \eqref{eq:tpc57-two-affine-forms} for which
\begin{enumerate}[label=\textup{(\roman*)},leftmargin=3em]
 \item \(\ell(t)\) belongs to the declared source-prime row support;
 \item the pair \((d\ell(t),r)\), the orbit time \(j\), and the opposite
       row \(\gamma\) satisfy every retained literal carrier condition;
 \item \(p(t)\) is in the declared terminal-prime support and its shaped
       terminal amplitude is nonzero.
\end{enumerate}
For terminal-admissible occupancy one does not additionally delete a zero
source coefficient; for physical occupancy one intersects this set with
the exact nonzero coefficient support.  With that convention, the fixed
divisor occupancy is the identity
\begin{equation}
 N_{\gamma,s,j,d}
 =\#\mathfrak T_{\gamma,s,j,d}
 =\sum_{t\in\mathbb Z}
   \mathbf 1_{\mathfrak T_{\gamma,s,j,d}}(t).
 \label{eq:tpc57-exact-ladder-occupancy}
\end{equation}
This definition retains the literal carrier and does not replace it by
all prime values of the two affine forms.

\begin{corollary}[Deterministic ladder capacity]
\label{cor:tpc57-ladder-capacity}
For source-prime support inside \([L,2L]\),
\begin{equation}
 N_{\gamma,s,j,d}
 \le1+\left\lfloor\frac{Lg}{r}\right\rfloor
 \le1+\left\lfloor\frac{LH}{r}\right\rfloor.
 \label{eq:tpc57-exact-ladder-capacity}
\end{equation}
In particular, \(r>LH\) forces \(N_{\gamma,s,j,d}\le1\).
\end{corollary}

\begin{proof}
The first affine form in \eqref{eq:tpc57-two-affine-forms} has spacing
\(r/g\), and \([L,2L]\) has diameter \(L\).  The second inequality uses
\eqref{eq:tpc57-literal-ladder-gcd}.  Every further literal condition
only deletes parameters.
\end{proof}

\subsection{What prime input remains}

After the deterministic reductions, actual occupancy on one fixed
divisor is a simultaneous-primality problem for the two nonproportional
affine forms
\begin{equation}
 \ell_0+\frac r g t,
 \qquad
 p_0+\frac{dj}{g}t,
 \label{eq:tpc57-two-prime-problem}
\end{equation}
whose determinant is the fixed nonzero integer
\(-h_0/g\).  The determinant rules out proportional forms and exposes the
local congruence data.  It does not provide a lower bound for prime
values.  In particular:
\begin{enumerate}[label=\textup{(\roman*)},leftmargin=3em]
 \item a lower bound on an individual growing ladder is a two-linear-form
       prime problem and is not supplied by the elementary geometry;
 \item an upper-bound sieve may thin the ambient ladder by logarithmic
       factors, but it neither proves terminal activation nor compares a
       collision second moment to the physical terminal energy;
 \item averaging over \(d,r,j,s\) may be more tractable than one fixed
       ladder, but uniformity in the moving coefficients, literal masks,
       residue classes, and physical coefficient weights requires a
       separate analytic theorem.
\end{enumerate}
Thus \cref{prop:tpc57-two-affine-reduction} is an exact arithmetic
reduction of the next gate.  It is not a singleton-density theorem, a
parity-breaking estimate, or a fixed-\(h_0\) prime-pair lower bound.
